% =============================================================================
%  Part V -- Axioms of Causal Physics
% =============================================================================

This part states the foundations of Causal Physics in two groups. The \emph{methodological
principles} say how the theory may be built and tested; they are rules of inquiry and say nothing
about nature. The \emph{physical axioms} say what the world is made of and how it behaves; they are
the content of the theory. For each physical axiom we state its status: whether it is an
\emph{input} (assumed), a \emph{derived} statement, or \emph{tested} by the laboratory.

No background manifold, metric, coordinate system or spacetime dimension is assumed. All such
structures, where they exist, are read from the order.

% =============================================================================
\section{Methodological principles}\label{sec:method-ax}
% =============================================================================

\begin{definition}[M1: Causal primacy]
No geometry, coordinate system, external clock or unit enters the causal layer. Geometry is a
description of the order, read from it by declared instruments.
\end{definition}

\begin{definition}[M2: Separation of layers]
Causal, bridge and observational quantities are kept apart (Sec.~\ref{sec:layers}). Observational
quantities never act on the causal dynamics, and bridge parameters never change its rules.
\end{definition}

\begin{definition}[M3: Order and number]
Every rule of the dynamics uses only the order and counts of events, or states of events derived
from them. Coordinates may appear in an instrument as a declared readout, never in a rule.
\end{definition}

\begin{definition}[M4: Causality and covariance of rules]
The state of an event depends only on its causal past. A rule gives the same result in every
inertial frame and in every order of computation compatible with the causal order (every linear
extension).
\end{definition}

\begin{definition}[M5: Finiteness]
A rule and its results are finite and well defined without a wall of the region or a regulator. A
result that holds only because of the boundary of a simulation is a property of the simulation.
\end{definition}

\begin{definition}[M6: Input and output]
Every statement is classified as an input (postulate, tuned parameter) or an output (it followed
without being put in). A test of an output that was put in as an input proves nothing.
\end{definition}

\begin{definition}[M7: Pre-registration and falsifiability]
Every test states its question, prediction, null model and numerical criterion before it is run,
and declares the level of conclusion it can support: a \emph{calibration} establishes an instrument and
says nothing about physics; a \emph{mechanism} shows that a mechanism works in a declared model; a
\emph{chain} shows that connected mechanisms reproduce a predicted behaviour; a \emph{prediction} is an
independent result fixed in advance. No output may claim more than its level allows. A changed criterion
makes a new revision of the test, and the earlier verdict stays. A new hypothesis may follow a failure, but
it may not change the criteria by which the earlier test failed. Negative results are kept.
\end{definition}

% =============================================================================
\section{Physical axioms}\label{sec:phys-ax}
% =============================================================================

\subsection{The substrate}

\begin{axiom}[Primacy of events]\label{ax:events}
Physical reality consists of elementary events. An event is an irreducible act of realisation, not a
persisting entity. It has no intrinsic position, duration, mass, energy, charge or scale. All
physical attributes are relational: they are patterns of causal connectivity among events.
\end{axiom}
\begin{status} Input. \end{status}

\begin{axiom}[Causal order]\label{ax:order}
Causal relations form a partial order: they are transitive and acyclic, and every interval contains
finitely many events. Closed causal loops do not exist. Time is the order itself; there is no
external time and no global simultaneity.
\end{axiom}
\begin{status} Input. The structure is a causal set. \end{status}

\begin{axiom}[Lorentz-invariant substrate]\label{ax:substrate}
On the plateau the order is, statistically, a uniform sprinkling of a Lorentzian manifold at density
$\rho$: the number of events in a region is $\rho$ times its proper volume, and no frame is
preferred. Dimension, distance and duration are read from order and number alone.
\end{axiom}
\begin{status} Input, supported by tests: the degree of an event is not bounded (F0.3), local
growth rules with bounded degree produce no manifold-like order (F6.1, F6.2), and geometry is read
blindly from the order (K.1, K.3, K.16b, K.16c). \end{status}

\subsection{Transport and matter}

\begin{axiom}[Causal continuity]\label{ax:continuity}
Influence is transported along the relations of the order and is conserved at every event:
\begin{equation}
  \sum_{\text{in}}F+s_v=\sum_{\text{out}}F+\Delta Q_v .
\end{equation}
What enters an event leaves it or is stored there.
\end{axiom}
\begin{status} Input. Tested consequence: Gauss's law for a steady source, $J/S=1.003$--$1.004$ in
2+1 and 3+1 (F1.1). \end{status}

\begin{axiom}[Finite capacity]\label{ax:capacity}
The flow an event can pass on is bounded by a capacity $C$. Capacity limits the flow through an
event, not the number of its relations.
\end{axiom}
\begin{status} Input. Supported by F1.8 (not yet robust to the seed). \end{status}

\begin{axiom}[Flows carry momentum]\label{ax:momentum}
A flow carries its four-velocity $u^a$ as a state passed from event to event. When flows merge, the
merged flow is the length of the summed four-momentum, $F=|\sum P|$, so that four-momentum is
conserved. The flows then have an energy--momentum tensor $T^{ab}$, Eq.~\eqref{eq:Tab}.
\end{axiom}
\begin{status} Input. Tested consequences: a flow with a carried direction is localised and moves with
its source, i.e.\ has inertia (F3.6, replicated); $T^{ab}$ is measured correctly for dust,
counter-streaming beams and an isotropic gas, and the merging rule conserves four-momentum exactly
(K.15, replicated). \end{status}

\subsection{Response and action}

\begin{axiom}[Local response]\label{ax:response}
Every event has a response variable $\Pi_v$, $0<\Pi_v\le1$, read as its clock rate, and a stiffness
$I_v\ge1$; its impedance is $Z_v=I_v/\Pi_v$. The axiom does not fix how $\Pi_v$ depends on the
load, on the surrounding order or on other causal variables; that dependence is a separate,
independently testable law of response.
\end{axiom}
\begin{status} Input. The law of response is not part of the axiom. In the linear regime it is
fixed by the dynamics of the order (Axiom~\ref{ax:dynamics}, Part~\ref{part:geometric}). \end{status}

\begin{axiom}[Action of a flow]\label{ax:action}
The action of a flow is its proper time counted by the order,
\begin{equation}
  S_m=-\sum m\,\Delta\tau ,
\end{equation}
where $\Delta\tau$ is the proper time between successive events of the flow, as given by the length
of the longest chain or by the number of elements of a small interval. The Causal Resistance
Functional $\mathcal R=\sum_v Z_v\Delta\tau$ is its impedance-weighted form.
\end{axiom}
\begin{status} Input. Derived consequences: on a fixed order the motion is geodesic in the conformal
geometry $\Pi^2\eta$ (tested: F2.2, F2.4, F2.5); the coupling of a flow to the order is
$\delta S_m/\delta h_{ab}=\tfrac12T^{ab}$, i.e.\ through the rest frame of the flow, and contributions
of different flows add (Sec.~\ref{sec:action}). How the weight $Z_v$ relates to the response of the
order beyond the linear regime is open. \end{status}

\subsection{Dynamics of the order}

\begin{axiom}[The order responds to matter]\label{ax:dynamics}
The order is dynamical. Its configurations are weighted by an action of the order, and matter
changes the order: near flows, causal cones narrow and events are added. Causality itself is
preserved (Axioms~\ref{ax:order} and~\ref{ax:continuity}). The candidate action of the order is the
Benincasa--Dowker action $S_\BD$ (Sec.~\ref{sec:bd}).
\end{axiom}
\begin{status} Input (that the order is dynamical and which action weights it). Derived and computed
consequences:
\begin{itemize}
\item the expected BD action has the Einstein--Hilbert stiffness in the static conformal sector,
  $6.0003$ against $6$ (F6.8a), in the mixed sector, $7.983$ against $8$ (F6.8b2), and in the optical sector,
  $2.002$ against $2$ (F6.8c2); with Axiom~\ref{ax:action} the linear response of the order is then
  the coupling of linearised general relativity to $T_{ab}$ (Part~\ref{part:geometric});
\item in 3+1 this response is the conservation of causal throughput, $\rho_c\,c_{\rm eff}=\text{const}$,
  which gives $\gamma=1$ (tested on discrete orders: F7.1c, raw $\gamma=1.000\pm0.045$, formal verdict
  undecided);
\item a substrate rule realising it makes pressure and motion gravitate and drags the cones
  (tested: F7.4a).
\end{itemize}
How the order grows (a sequential growth law or a sum over orders weighted by $S_\BD$) is
open. \end{status}

% =============================================================================
\section{What is derived, and from what}\label{sec:ax-summary}
% =============================================================================

\begin{center}
\begin{tabularx}{\linewidth}{>{\raggedright\arraybackslash}X>{\raggedright\arraybackslash}p{2.6cm}>{\raggedright\arraybackslash}p{4.8cm}}
\toprule
statement & from & status\\
\midrule
time as order, no global simultaneity & Ax.~\ref{ax:order} & derived\\
geometry read from order and number & Ax.~\ref{ax:substrate} & tested (K.16b/c)\\
Gauss's law & Ax.~\ref{ax:continuity} & tested (F1.1)\\
inertia of localised matter & Ax.~\ref{ax:momentum} & tested (F3.6)\\
conservation of four-momentum in merging & Ax.~\ref{ax:momentum} & tested (K.15)\\
geodesic motion & Ax.~\ref{ax:action} & tested (F2)\\
coupling through $T_{ab}$ in the rest frame of each flow & Ax.~\ref{ax:action} & derived\\
Einstein--Hilbert stiffness of the order & Ax.~\ref{ax:dynamics} & computed, all three sectors (F6.8a, F6.8b2, F6.8c2)\\
throughput principle, $\gamma=1$ in 3+1 & Ax.~\ref{ax:action}, \ref{ax:dynamics} & derived; F7.1c\\
tensor coupling, frame dragging & Ax.~\ref{ax:action}, \ref{ax:dynamics} & derived; F7.4a\\
non-linear completion by general relativity & Ax.~\ref{ax:dynamics} + uniqueness & derived (programme)\\
\bottomrule
\end{tabularx}
\end{center}
